\documentclass[11pt]{article}
\usepackage{amsfonts,amsmath,amssymb,bm,booktabs,enumitem,array}
\usepackage[letterpaper,textwidth=6.3in,textheight=8.7in]{geometry}
\usepackage[T1]{fontenc}
\usepackage{xcolor}
\usepackage{listings}
\usepackage[authoryear,round]{natbib}
\usepackage[breaklinks=true,colorlinks,linkcolor=blue,citecolor=blue,urlcolor=blue]{hyperref}

\lstset{basicstyle=\ttfamily\small,columns=fullflexible,keepspaces=true,frame=single,framerule=0.3pt,rulecolor=\color{black!40},xleftmargin=0.5em,framexleftmargin=0.5em,breaklines=true,showstringspaces=false,aboveskip=0.6em,belowskip=0.6em}
\newcommand{\code}[1]{\texttt{#1}}
\newcommand{\E}{{\mathbb E}}
\newcommand{\PP}{{\mathbb P}}
\newcommand{\iu}{{\mathrm i}}
\newcommand{\dr}[1]{{\mathrm d}#1}
\newcommand{\Rea}{\operatorname{Re}}
\newcommand{\phih}{\hat\phi}
\newcommand{\tauh}{\hat\tau}
\newcommand{\Fk}{\textsf{DK}}
\newcommand{\Fu}{\textsf{DU}}
\newcommand{\Ft}{\textsf{DT}}
\newcommand{\Sk}{\textsf{SK}}
\newcommand{\Su}{\textsf{SU}}
\newcommand{\St}{\textsf{ST}}
\setlist[description]{leftmargin=1.5em,itemsep=0.35em,style=nextline,font=\normalfont\ttfamily}
\setlength{\parskip}{0.25em}

\title{The Exact Limiting Distribution of the Dickey--Fuller $t$-Ratio\\[0.6ex]\large MATLAB programs: user manual}
\author{Yong Bao\\\normalsize Purdue University \and Simon A. Broda\\\normalsize Lucerne University of Applied Sciences \and Raymond Kan\\\normalsize University of Toronto and\\\normalsize Hong Kong University of Science and Technology}
\date{Version 1.0, November 2026}

\begin{document}
\maketitle

\begin{abstract}
\noindent This manual documents the MATLAB programs that accompany \citet*{BaoBrodaKan2026l} (``the paper'').  The programs compute the limiting distribution function, density, quantiles and moments of the least squares $t$-ratio $\tauh_0=(\phih-\phi_0)/\operatorname{se}(\phih)$ in a Gaussian first-order autoregression under the local alternatives $\phi=1+c/n$, for the regression with no deterministic terms, an intercept, or an intercept and a linear trend, for the null value $\phi_0=1+c_0/n$ ($c_0=0$: the Dickey--Fuller tests), and with a zero, a fixed or a stationary initial observation of the order of $\sqrt n$.  Two exact representations are implemented, the line integral of Theorem~2 of the paper and the conditional integral of Theorem~3, and they check each other.  Section~\ref{sec:start} is a quick start, Section~\ref{sec:user} describes the user-level routines, Section~\ref{sec:engines} the two engines, Section~\ref{sec:accuracy} accuracy and cost, Section~\ref{sec:validation} the validation suite and the reproduction of the paper, and Section~\ref{sec:files} lists the files.
\end{abstract}

\tableofcontents

\section{Quick start}
\label{sec:start}

Add the folder \code{public} to the MATLAB path.  The model is
\[
y_t=\alpha+\delta t+\phi y_{t-1}+\sigma e_t,\qquad t=1,\ldots,n,\qquad e_t\sim N(0,1)\text{ i.i.d.},
\]
with $\phi=1+c/n$, no drift, and $y_0=\sigma\sqrt n\,\gamma_0$.  The case code names the start (\textsf{D}eterministic, that is fixed, or \textsf{S}tationary) and the regressors (intercept \textsf{K}nown to be zero; \textsf{U}nknown intercept; intercept and \textsf{T}rend):
\begin{center}
\begin{tabular}{@{}lll@{}}
\toprule
Code & Regression & Start\\
\midrule
\code{'DK'} & $y_t$ on $y_{t-1}$ ($\hat\tau$) & $\gamma_0$ fixed (\code{opts.g0}, default $0$)\\
\code{'DU'} & $y_t$ on $1$, $y_{t-1}$ ($\hat\tau_\mu$) & $\gamma_0$ fixed\\
\code{'DT'} & $y_t$ on $1$, $t$, $y_{t-1}$ ($\hat\tau_\tau$) & $\gamma_0$ fixed\\
\code{'SK'}, \code{'SU'}, \code{'ST'} & the same & stationary, $J(0)\sim N(0,-1/(2c))$, $c<0$\\
\bottomrule
\end{tabular}
\end{center}
All routines take \code{(\ldots,cas,c,opts)}: the case, the local-to-unity parameter $c$ (default $0$, the unit root) and an optional structure of options, of which \code{c0} (the local null value, default $0$), \code{g0} (the fixed start, default $0$) and \code{method} (\code{'auto'}, the default, \code{'single'} or \code{'double'}) are the ones a user needs.
\begin{lstlisting}
ltratio_icdf([0.01 0.05 0.10], 'DU')        % asymptotic critical values of tau_mu
ltratio_cdf(-2.86157, 'DU', -5)              % limiting power of the 5% test at c = -5
ltratio_cdf(-2.86157, 'SU', -5)              % the same with a stationary start
ltratio_cdf(-2.86157, 'DU', -5, struct('g0', 1.5))   % fixed start y_0 = 1.5 sigma sqrt(n)
ltratio_cdf(0, 'DK', 2, struct('c0', 2))     % t-ratio centered at the true value, c = 2
ltratio_pdf(-1:0.5:1, 'DK')                  % the density of tau at five points
ltratio_mom(4, 'DT')                         % E tau_tau^j, j = 1..4
\end{lstlisting}
The first call returns $-3.43035$, $-2.86157$ and $-2.56679$; the second $0.1174$.  A value of the distribution function or the density takes a fraction of a second, a quantile a few values, and a set of moments about a second.

\section{The user-level routines}
\label{sec:user}

\paragraph{Common arguments.}  \code{cas} is the case code (in either letter case), \code{c} a finite real scalar (negative for a stationary start), and \code{opts} a structure.  The arguments are checked by \code{ltratio\_args}, which rejects unknown codes, a stationary start with $c\ge0$, a fixed start \code{g0} in a stationary case and an unknown \code{method}.  Fields of \code{opts} other than \code{c0}, \code{g0} and \code{method} are passed to the engine in use (Section~\ref{sec:engines}).  Vector arguments are evaluated elementwise.

\begin{description}
\item[{[F,info] = ltratio\_cdf(x,cas,c,opts)}]
The limit of $\PP[\tauh_0\le x]$.  At \code{x = -Inf}, \code{Inf} and \code{NaN} the result is $0$, $1$ and \code{NaN}.  \code{info.method} is the representation used and \code{info.time} the elapsed time; with \code{'double'} the diagnostics of \code{tr\_limdist} are added (\code{info.status}, \code{info.err}, \ldots; a nonzero status raises the warning \code{tr\_limdist:status}).
\item[{[f,info] = ltratio\_pdf(x,cas,c,opts)}]
The limiting density; arguments and options as for \code{ltratio\_cdf}.  In case \Fk\ with $c_0>0$ the density is infinite at $x=-\sqrt{2c_0}$ (a logarithmic singularity, where the two zeros of the kernel's edge merge) and \code{Inf} is returned there.
\item[{[q,info] = ltratio\_icdf(p,cas,c,opts)}]
The limiting $p$-quantiles, $p\in(0,1)$, by Newton's method on \code{ltratio\_cdf} with \code{ltratio\_pdf} as derivative, safeguarded by bisection, started at the normal approximation from the first two moments.  \code{opts.tol} ($10^{-10}$) is the tolerance on $q$; \code{info.F} is the distribution function at $q$, \code{info.iter} the iteration count and \code{info.conv} the convergence flag (a failure raises \code{ltratio:icdf}).
\item[{[m,info] = ltratio\_mom(p,cas,c,opts)}]
The moments $\E(\tau^j)$, $j=1,\ldots,p$, of the limit $\tau=(U-c_0V)/\sqrt V$, as one-dimensional integrals of the joint moment generating function at real arguments (Proposition~3 of the paper), through \code{tr\_limmom}; the options \code{c0} and \code{g0} apply, \code{method} does not.  \code{info} holds the error estimate of the integral per moment and the time.
\end{description}

\paragraph{The two representations.}  With $U=\int_0^1\tilde J\dr J$ and $V=\int_0^1\tilde J^2\dr r$, $J$ the Ornstein--Uhlenbeck process $\dr J=cJ\dr r+\dr W$ started at $\gamma_0$ and $\tilde J$ its residual on the deterministic terms, the limit is $\tau=(U-c_0V)/\sqrt V$, and the joint moment generating function of $(U,V)$ is $e^{-s/2}D(s,t)^{-1/2}$ with $D$ a polynomial of degree one (\Fk) or two in $s$ (Proposition~1 of the paper), to whose logarithm a fixed start adds $\gamma_0^2Q(s,t)$ in closed form (Proposition~2).  The \emph{line integral representation} (\code{method = 'single'} or \code{'line'}, Theorem~2) inverts in $s$ explicitly, which leaves a scaled $\chi^2_1$ density (\Fk) or a Bessel $K_0$ density of $U+\frac12$ under the tilt $e^{tV}$, and integrates its Gaussian transforms along the line $\Rea t=\tau$; its engine is \code{dfx\_core}.  The \emph{conditional integral representation} (\code{method = 'double'} or \code{'conditional'}, Theorem~3) keeps the outer integral over the value $v$ of $V$ and inverts in $t$, for each $v$, along a path through a saddle point; its engine is \code{tr\_limdist}.  A fixed start $\gamma_0\ne0$ in cases \Fu\ and \Ft\ makes the tilted density noncentral and without a closed form, and only the conditional integral is available there; \code{ltratio\_cdf} and \code{ltratio\_pdf} then use it whatever \code{method} says.  Everywhere else both are available and agree to about $10^{-12}$ in the distribution function (\code{ltratio\_test}, Section~\ref{sec:validation}).  With \code{method = 'auto'}, the default, \code{ltratio\_method} chooses point by point the representation that is the faster one in these programs, by their own timings (\code{validation/method\_sweep.m}): the line integral at $c_0=0$, where it is $1.2$ to $60$ times faster; with $c_0\ne0$ the conditional integral in cases \Fk, \Fu, \Ft\ and \Sk, except for $c_0>0$ at $x\ge0$, where the line integral has no breakpoint and is $3$ to $30$ times faster (with $c_0<0$ its saddle factor costs up to $11$ seconds at $x$ of $4$ to $8$, while the conditional integral stays below half a second); and the line integral in cases \Su\ and \St, whose conditional integral costs half a second a value, except for $c_0<0$ at $x\ge0$.  The C and Julia implementations apply the same rule for every start; the exception for cases \Su\ and \St\ at $x<0$ is the one place where the MATLAB timings differ.  A point at which the conditional integral reports that its error estimate is not met (\code{info.status}) gets the line integral's value instead, in all three implementations (\code{info.nfallback} counts them); with \code{method = 'double'} the flagged value is returned, with a warning.

\section{The two engines}
\label{sec:engines}

\paragraph{The line integral: \code{dfx\_core}.}  Called through \code{dfx\_cdf(q,model,c,opts)} and the front ends \code{dfx\_pdf}, \code{dfx\_quantile}, \code{dfx\_moments}, \code{dfx\_tstar}, with the model number $1$, $2$, $3$ for \textsf{K}, \textsf{U}, \textsf{T} and \code{opts.start} equal to \code{'zero'} (default) or \code{'stationary'}; \code{ltratio\_dist} does the mapping.  The port follows the Julia implementation \code{DFExact.jl} of the line integral function by function.  The line is $\Rea t=\tau$ with $\tau=1$ for $c\le0$ and a zero start and $\min(1,t_*/2)$ otherwise, $t_*$ the convergence abscissa of the moment generating function of $V$ (\code{dfx\_tstar}; located from the cancellation-free Girsanov forms for $c>2$, so that $c$ up to $20$ and beyond is reached).  The integrand is split, by an entire partition of unity inside the inner transforms, into a smooth part, integrated by the trapezoidal rule in $w=\tau\sinh u$ with the step $0.075$ up to $w=4900$, and oscillating parts $G_j(w)e^{-\iu wv_j^2}$, one for each positive root $v_j$ of $X(v)=c_0v^2+xv+\frac12$ (one root, $1/(2|x|)$, at $c_0=0$ and $x<0$), integrated by Filon's method on geometric panels with $20$ Gauss--Legendre nodes; for $c>0$, and in the upper tail with $c_0<0$, the smooth part carries a saddle-point factor whose phase grows like $x^2$, and it is integrated on Gauss--Legendre panels whose node counts follow the phase.  The inner transforms are Gauss--Legendre integrals on complex paths on which the Gaussian kernel decays; with $c_0\ne0$ the paths between the breakpoints (the roots of $X$, and the vertex of the parabola for $c_0>0$) are straight pieces along the directions of steepest descent, joined through the real axis where a direct connection would cross a region in which the Gaussian grows (Appendix~B.1 of the paper).  Options (for checking): \code{tau} (the line), \code{hu} ($0.075$), \code{N} ($20$), \code{pfac} (subdivides the Filon panels), \code{glfac} (multiplies the node counts of the Gauss--Legendre route), \code{wmax} ($4900$), \code{gen} (\code{true}: the general kernel also at $c_0=0$), \code{uniform} (the trapezoidal rule for the whole integrand); \code{dfx\_core('nodefac',2)} doubles every node count until reset with \code{1}.  The moments of \code{dfx\_moments} use the form of Proposition~1 and lose accuracy for $c>2$; \code{ltratio\_mom} uses \code{tr\_limmom}.

\paragraph{The conditional integral: \code{tr\_limdist}.}  \code{[F,info] = tr\_limdist(x,cas,c,opts,mode)}, \code{mode} \code{'cdf'} or \code{'pdf'}, with \code{opts.c0} (default $c$ there, not $0$: \code{ltratio\_dist} passes \code{c0} explicitly) and \code{opts.g0}.  In case \Fk\ it uses part (i) of Theorem~3, whose inner probability is an error function (with a fixed start, the probability $\PP[X^2\le2a+\gamma_0^2]$ of a noncentral normal); with a zero start in cases \Fu, \Ft\ and with a stationary start in all three cases part (ii), whose integrands are in closed form; and in cases \Fu, \Ft\ with a fixed start part (i) with the noncentral probability (\code{Pnc}), a one-dimensional integral with \code{Nl} ($96$) Gauss--Legendre nodes, so a triple integral in all.  With a zero start the integrands are the Girsanov forms (\code{iipar}), analytic beyond $t_*$, and the inner inversions run through the saddle point of the whole integrand on two paths (a parabola turning into a wedge; a vertical line in the $\vartheta$-plane; in case \Fk\ a third route, $\operatorname{erf}=1-\operatorname{erfc}$, inverts the density of $V$ and the remainder separately), with $96$ trapezoidal nodes after a sinh substitution, doubled up to three times until the nested-rule estimate is below $10^{-13}$; a value is accepted when two paths agree to that tolerance, or when the other paths are consistent with the best one within ten times the sums of the estimates (the nested rule alone being open to aliasing, \code{agreepaths}); otherwise the node counts are doubled up to three more times, after which the discrepancy is reported as the error.  With an initial observation the path is the Talbot parabola through the saddle point of $\log M_V(t)-tv$, with the sign conditions of Theorem~3 checked at every node (\code{info.nbad} counts violations) and the branch of $\log D$ continued along it; \code{Nv} ($96$) nodes, doubled up to twice where the nested estimate exceeds $10^{-9}$ of the largest column of the piece and $10^{-11}$ divided by the range of $v$ (\code{info.nunres} counts the columns still unresolved).  The initial-value terms come from \code{tr\_nsstart} (Proposition~2: $Q(s,t)$ in closed form, from Nabeya and S\o rensen's determinants in cases \Fu, \Ft\ and the Riccati solution \code{tr\_jbstart} in case \Fk); the squared means of the noncentral form are read off $Q$ at three values of $s$.  The outer integral over $v$ is a cosine-mapped Gauss--Legendre rule in $\log v$ over $[e^{-8},e^6]\,\E V$ with \code{M} ($64$; $32$ in case \Fk\ with a zero start) nodes per part of length at most \code{plen} ($2$), breakpoints at the zeros of $a(v)=\frac12+x\sqrt v+c_0v$, grading toward the step of the integrand in the far upper tail for $c>c_0$ (width below \code{stepw}, $0.1$) and bisection, up to \code{maxlev} ($12$) times, where the Legendre-coefficient estimate exceeds \code{otol} ($10^{-12}$).  \code{info.err} is the error estimate per point and \code{info.status} is $0$ when clean, otherwise $1$ (sign violations) plus $2$ (unresolved inner integrals, or \code{info.err} above \code{errtol}, $10^{-10}$); a nonzero status raises the warning \code{tr\_limdist:status}.  Further options: \code{blo}, \code{bhi} (the range in $\log v$ relative to $\log\E V$), \code{region} (checks the sign conditions on that many intermediate paths), \code{oldDK}, \code{oldII} (the earlier forms through the saddle point of $M_V$ alone, which fail for $c$ above about $2$ and $5$).

\paragraph{The moment generating function.}  \code{tr\_cflim(t1,t2,cas,c,form)} evaluates the Nabeya--Tanaka determinants of Proposition~1 at complex arguments (31-term Taylor series in $z$ for $|z|<4$ from \code{nt\_series}, closed forms in $\vartheta=(-z)^{1/2}$ and $e^{-2\vartheta}$ beyond), with the branch of $\log D$ followed along two legs and the roots of $D$ in $s$ (the weights $\kappa_k$ of $U+\frac12$ under the tilt); the forms \code{'F'}, \code{'B'}, \code{'BC'} and \code{'ApB'} return the coefficients, the weights, the branch-continued logarithm along a path, and the moment generating function of $U$ and $V$ together.  \code{tr\_limmom(p,cas,c,opts)} computes the moments by Cauchy integrals on a circle in the argument of $U$ and the trapezoidal rule in $\log s$, from the Girsanov forms with a zero start and with the initial-value terms of \code{tr\_nsstart} otherwise.

\section{Accuracy and cost}
\label{sec:accuracy}

Both representations are accurate to about $10^{-11}$ in the distribution function and $10^{-10}$ relative in the density over $c\in[-30,20]$, $c_0\in[-40,10]$ and the six cases, and to $10^{-11}$ with a fixed start in cases \Fu, \Ft\ for $c\le2$ (Section~7 and Appendix~B of the paper).  The error of the line integral is absolute rather than relative: its terms are of order one, and it carries an error of $10^{-14}$ to $10^{-21}$ whatever the size of the value, which finer rules do not reduce ($2\times10^{-14}$ in the lower tails of cases \Ft\ and \St\ with $c_0\ne0$).  When \code{opts.method} is not given, \code{ltratio\_pdf} and \code{ltratio\_cdf} therefore recompute by the conditional integral every value of the line integral below $10^{-3}$ (a density, or a distribution function in the lower tail), with its error estimate in \code{info.err}.  The relative accuracy then holds in the tails as well, down to that estimate; far out, below the error of either method, a computed density can be a tiny negative number, and it is returned as zero.  Doubling every node count of the line integral changes the distribution function by at most $4\times10^{-13}$; the two representations agree to $4\times10^{-12}$ at worst (Table~1 of the paper, reproduced by \code{validation/acc\_table.m}).  With a fixed start in cases \Fu, \Ft\ at $c$ above about $2.5$ the Talbot path of the conditional integral no longer satisfies the sign conditions everywhere, and \code{tr\_limdist} warns; the line integral does not cover that configuration.  A value of the distribution function or the density takes, in MATLAB on one thread of a laptop, $0.1$ to $0.3$ seconds by the line integral for $c_0=0$ and $0.3$ to $5$ seconds with $c_0\ne0$ (the upper end in the upper tail at $c=5$, $c_0=-5$), and $0.05$ to $0.4$ seconds by the conditional integral with a zero start and $0.3$ to $1.3$ seconds with an initial observation; a quantile takes a few values and a set of four moments about a second.  The Julia implementation of the line integral, \code{DFExact.jl}, takes $7$ to $17$ milliseconds per value on six threads.

\section{Validation and reproduction of the paper}
\label{sec:validation}

\code{ltratio\_test} runs a hundred and fourteen checks in about two minutes and reports a failure count: the two representations against each other at the mean plus $-2.5$ to $2.5$ standard deviations in the six cases, at $c\in\{0,-5,-10\}$, $c=2$ and $5$ and $c_0\ne0$; the exact values $\PP[\hat\tau\le0]=\operatorname{erf}\{(2w)^{-1/2}\}$, $w=(e^{2c}-1)/(2c)$, for $c\in\{-10,0,2,5,20\}$, $\PP[\hat\tau_\mu\le0]=0.95756559$ from the coefficient statistic, and $\PP[U-c_0V\le0]$ at $c=10$, $c_0=5$ from an independent one-dimensional inversion; the distribution function at the asymptotic critical values of \citet{MacKinnon2010} and the quantile function; the initial observation (continuity as $\gamma_0\to0$; the stationary start as the 24-point Gauss--Hermite average of the fixed start; the noncentral probability $\PP[J(1)^2\le1+\gamma_0^2]$ in case \Fk; the rank-two structure of the initial-value term); the moments against \citet{Nabeya1999}, against \code{dfx\_moments} and against the integral of the density; the density against fifth-order differences of the distribution function; and the argument checks.  \code{ltratio\_paper} recomputes the tables of the paper and prints each value next to the published one (about ten minutes); \code{validation/acc\_table.m} produces the accuracy table and its LaTeX rows.

\section{Files}
\label{sec:files}

\begin{description}
\item[ltratio\_cdf.m, ltratio\_pdf.m, ltratio\_icdf.m, ltratio\_mom.m] the user-level routines; \code{ltratio\_dist.m} (dispatch to the engines), \code{ltratio\_method.m} (the automatic choice of the representation) and \code{ltratio\_args.m} (argument checks).
\item[dfx\_core.m, dfx\_cdf.m, dfx\_pdf.m, dfx\_quantile.m] the line integral: the engine and its front ends;
\item[dfx\_moments.m, dfx\_tstar.m, dfx\_series.m, dfx\_test.m] its moments, convergence abscissa, series coefficients and test.
\item[tr\_limdist.m, tr\_limmom.m] the conditional integral and the moments;
\item[tr\_cflim.m, tr\_nsstart.m, tr\_jbstart.m, nt\_series.m, gen\_series.py, faddeeva.m] the moment generating function, the initial-value terms, the series coefficients and the Faddeeva function.
\item[ltratio\_test.m, ltratio\_paper.m, validation/] validation and reproduction.
\item[Contents.m, README.md, LICENSE, doc/] this manual.
\end{description}

\begin{thebibliography}{9}
\bibitem[Bao et~al.(2026a)]{BaoBrodaKan2026l}
Bao, Y., Broda, S.~A. and Kan, R. (2026a).  The exact limiting distribution of the Dickey--Fuller $t$-ratio.  Manuscript.
\bibitem[Bao et~al.(2026b)]{BaoBrodaKan2026}
Bao, Y., Broda, S.~A. and Kan, R. (2026b).  The exact finite-sample distribution of the Dickey--Fuller $t$-ratio.  Manuscript.
\bibitem[MacKinnon(2010)]{MacKinnon2010}
MacKinnon, J.~G. (2010).  Critical values for cointegration tests.  Queen's Economics Department Working Paper 1227.
\bibitem[Nabeya(1999)]{Nabeya1999}
Nabeya, S. (1999).  Asymptotic moments of some unit root test statistics in the null case.  \emph{Econometric Theory} 15, 139--149.
\end{thebibliography}

\end{document}
